% SOURCE: MIT 18.335J Spring 2019, Lecture 24 sparse matrix handout
% (Per-Olof Persson slides dated November 28, 2006), and Week 8 L23 summary.
% ADAPTATION: self-authored connected explanation, graph arguments, cost derivations and examples.
\originalchapter{稀疏存储、填充与直接法的复杂度}\label{ch:sparse}
稀疏矩阵并不只是一个有很多零的稠密矩阵. 我们必须同时改变存储、运算和消元顺序, 才能真正从零结构中获益.

\originalsection{稀疏结构与存储方式}
若一个 $n\times n$ 矩阵只有 $s$ 个非零元, 用稠密数组仍需 $n^2$ 个数的位置. 稀疏格式记录非零数值及其索引, 使存储约为 $O(s+n)$, 矩阵向量积约为 $O(s+n)$. 这里的 $n$ 包含输出向量初始化和行列指针, 不应在矩阵可能有大量空行时被省略.

结构矩阵与稀疏矩阵是不同概念. Toeplitz 矩阵可能几乎没有零, 但可用少量参数描述并快速相乘; 有些算子甚至不显式存储任何矩阵. Krylov 方法只需要快速的矩阵向量积, 因而适用范围比显式稀疏矩阵更广.

坐标格式 COO 存三个数组 $(i_\ell,j_\ell,a_\ell)$, 表示各项所在行、列和数值. 它便于组装, 特别是有限元计算中不同单元会向同一个矩阵位置贡献数值. 组装完成后应合并重复索引, 并按实际策略删除精确零; 否则存储中的条目数与数学上的非零元数不一致.

\begin{definition}
压缩行格式 CSR 用三个数组表示矩阵: 数值数组、列索引数组, 以及长度 $n+1$ 的行指针数组. 第 $i$ 行的所有条目位于第 $i$ 和第 $i+1$ 个行指针之间. 压缩列格式 CSC 交换行列角色, 按列保存数值与行索引.
\end{definition}
为避免编程语言下标差异, 下面例子用从 $1$ 开始的数值与索引数组, 并规定最后一个指针等于 $s+1$.

\begin{example}
把
\[
A=\begin{bmatrix}
4&0&-1&0\\
0&3&0&2\\
-1&0&5&0\\
0&2&0&6
\end{bmatrix}
\]
写成 CSR, 并计算其一次矩阵向量积的组织方式.
\end{example}
\begin{solution}
按每行列号递增排列, 数值数组为 $(4,-1,3,2,-1,5,2,6)$, 列索引为 $(1,3,2,4,1,3,2,4)$, 行指针为 $(1,3,5,7,9)$. 给定 $x$, 第 $i$ 个输出为该行对应条目的 $a_\ell x_{j_\ell}$ 之和. 此例共用 $8$ 次乘法, 每行两项合并需一次加法. 这些条目逐行连续读取, 不用扫描 $16$ 个稠密位置.
\end{solution}
CSR 适合逐行矩阵向量积, CSC 适合按列访问和许多稀疏分解. 这不意味着另一格式无法完成同一个数学运算, 而是局部访问、输出写入和库实现的代价不同. 若行内索引已经排序, 查找一个指定位置可二分搜索, 但插入新条目可能需要搬移后面的整个数组. 因此应区分\emph{组装格式}和\emph{计算格式}, 不宜在每次更新时都向压缩数组中插入一个元素.

稀疏矩阵相加或相乘也可能失去稀疏性. 若 $A_{ik}\ne0$ 且 $B_{kj}\ne0$, 则乘积的 $(i,j)$ 项有潜在贡献. 许多不同的 $k$ 会汇聚到同一个位置, 也可能恰好相消. 运算复杂度依赖这些结构交集和输出非零数, 不能只用输入各自的非零数简单相加来界定矩阵乘法.

\originalsection{消元为何产生填充}
对 Hermitian 正定矩阵, 先考虑消去一个主元的分块形式
\[
A=\begin{bmatrix}a&c^*\\c&B\end{bmatrix},\qquad a>0.
\]
Cholesky 或对称消元后, 其余变量面对的 Schur 补为
\begin{equation}\label{eq:sparse-schur}
S=B-\frac{cc^*}{a},\qquad S_{ij}=B_{ij}-\frac{c_i\overline{c_j}}a.
\end{equation}
若 $c_i,c_j\ne0$ 而 $B_{ij}=0$, 新值就是非零的 $-c_i\overline{c_j}/a$. 原来没有直接耦合的两个变量, 因为共同连接被消去的变量而出现了耦合. 这叫\emph{填充}.

\begin{definition}
对对称非零结构, 建立无向图 $G(A)$: 每个变量为一个顶点, 不同顶点 $i,j$ 间有边当且仅当结构上允许 $A_{ij}\ne0$. 消去一个顶点时, 把其尚未消去的邻居连接成团, 再删除该顶点. 新增边称为符号填充.
\end{definition}
式 \eqref{eq:sparse-schur} 就是这个图规则的代数依据. 若某条边原来已有数值, Schur 补更新可能出现精确相消; 符号分析通常按潜在非零处理, 忽略特殊数值相消, 因而给出预计结构或上界. 对一般数值, 这些潜在填充通常真的出现. 不能把每一条符号边无条件说成浮点输出中的严格非零.

\begin{example}\label{ex:sparse-star}
取
\[
A=\begin{bmatrix}
4&-1&-1&-1\\
-1&2&0&0\\
-1&0&2&0\\
-1&0&0&2
\end{bmatrix}.
\]
比较先消去中心变量 $1$, 与先消去叶子 $2,3,4$ 的填充.
\end{example}
\begin{solution}
叶子块为 $2I$, 中心对叶子块的 Schur 补为 $4-3/2=5/2>0$, 所以 $A$ SPD. 若先消去中心, 剩余矩阵为
\[
2I-\frac14\begin{bmatrix}1&1&1\\1&1&1\\1&1&1\end{bmatrix},
\]
三个叶子之间都出现 $-1/4$, 形成完整三角形. Cholesky 下三角因子一般有 $4+3+3=10$ 个非零条目, 即完整下三角.

若先消去任一叶子, 它只有一个剩余邻居, 不会产生邻居之间的新边. 三个叶子依次消去, 只把中心对角依次减少 $1/2$, 最后为 $5/2$. 下三角因子只需 $4$ 个对角条目和 $3$ 条星形边, 共 $7$ 个非零条目. 矩阵的数值和特征值没有改变, 只是变量顺序改变, 存储和后续更新就不同.
\end{solution}
若 $P$ 是置换矩阵, 对 $PAP^*$ 求分解并用相应顺序解方程, 得到的原解仍相同. 对 SPD 矩阵, 这样的对称置换保持正定性, 可以在数值分解前专门选择顺序以减少填充.

\originalsection{带宽、最小度与分隔子排序}
半带宽 $w$ 意味着 $A_{ij}=0$ 在 $|i-j|>w$ 时成立. 对 SPD 带状矩阵, 按自然顺序做 Cholesky 不会产生带外填充. 原因是消去第 $j$ 个变量时, 尚存邻居都在 $j+1,\ldots,j+w$ 中, 它们之间的距离不超过 $w$. 归纳即可保持带宽.

每列最多有 $w$ 个下三角非零元. 一次主元消去的稠密外积更新最多需要 $O(w^2)$ 工作, 所以带状 Cholesky 的上界为
\[
\text{存储 }O(nw),\qquad\text{分解 }O(nw^2),\qquad
\text{每个右端回代 }O(nw).
\]
三对角矩阵 $w=1$ 因此能在线性时间和线性存储中直接求解. 这正是前面 Lanczos 小三对角问题很便宜的一个原因.

带宽压缩方法如反向 Cuthill--McKee 试图把相邻顶点排到接近位置, 对长而窄的网格往往有效. 但较小带宽不等价于最小填充: 填充受局部邻居团结构和整个消元树影响, 不是只看最远的一条边.

最小度方法每一步消去当前剩余图中度数较小的顶点. 一个度数 $d$ 的顶点最多可能在其邻居间引入 $d(d-1)/2$ 条边, 因而低度数是合理的局部启发. 但当前度数会随填充改变, 而局部最小并不保证全局最佳. 实际的近似最小度方法通过压缩表示和度数估计降低维护图的成本.

\begin{definition}
图的分隔子是一个顶点集合 $S$, 删除 $S$ 后图分成彼此不相连的多个部分. 嵌套分割排序先递归排列这些部分, 再把分隔子 $S$ 的变量排在最后.
\end{definition}
把分隔子放在最后, 就能先分别消去子区域, 不让一个区域的内部消元直接把另一个区域的全部内部变量连接起来. 消元最终汇集到较小的界面. 对二维规则网格, 一条网格线可把区域近乎平分, 分隔子大小约为未知数总量的平方根; 对三维网格, 分隔面大小约为总量的 $2/3$ 次方. 二维和三维稀疏直接法的规模差异, 正是由这些界面大小引起.

求最优填充顺序本身是困难的组合优化问题. 实际排序常把嵌套分割与近似最小度混合使用. 不同图、机器内存、并行要求和数值主元策略下, 没有一个无需检验的万能顺序.

\originalsection{二维与三维网格的复杂度推导}
现在把 $N$ 用作未知数总数, 把 $m$ 用作每个方向的网格点数, 避免把二者混在一起. $d$ 维 $m\times\cdots\times m$ 网格有 $N=m^d$ 个未知数, 最近邻差分矩阵有 $O(N)$ 个非零元.

二维自然逐行编号的半带宽约为 $m=\sqrt N$. 带状估计给出存储 $O(N^{3/2})$, 分解 $O(N^2)$. 即使输入只有 $O(N)$ 个非零元, 填充仍可能比输入大得多. 对规则二维 Poisson 网格, 这些量级确实描述自然顺序的主要成本. 嵌套分割能显著改善它们.

为了精确说明估计的适用范围, 下述结论采用规则网格或具有相同比例分隔子的网格模型. 每次近乎平分区域, 在含 $N_s$ 个未知数的区域中, 假设其分隔子以及消元时的相关前沿大小均为 $O(N_s^{(d-1)/d})$. 最长边依次分割的规则网格满足这一量级. 消元前沿按稠密矩阵存储和更新, 一个大小 $s$ 的前沿需要 $O(s^2)$ 存储和 $O(s^3)$ 工作. 这些是本次分析的结构假设, 不是任意稀疏图自动拥有的性质.

\begin{proposition}[网格模型下的嵌套分割成本]\label{prop:nested-dissection-cost}
在上述模型中, 嵌套分割稀疏 Cholesky 的分解成本和因子存储量满足
\[
\begin{array}{c|ccc}
&d=1&d=2&d=3\\\hline
\text{因子存储}&O(N)&O(N\log N)&O(N^{4/3})\\
\text{分解工作}&O(N)&O(N^{3/2})&O(N^2)
\end{array}
\]
每个右端的前代、回代成本与因子非零数同阶.
\end{proposition}
\begin{proof}
在二叉递归的第 $\ell$ 层, 有 $O(2^\ell)$ 个区域, 每个含 $O(N/2^\ell)$ 个变量. 记分隔子大小
$s_\ell=O((N/2^\ell)^{(d-1)/d})$. 层存储上界为 $O(2^\ell s_\ell^2)$, 层工作为 $O(2^\ell s_\ell^3)$, 直到区域大小降为常数, 共 $O(\log N)$ 层.

当 $d=2$, 每层存储为 $O(N)$, 所以总量 $O(N\log N)$. 层工作为 $O(N^{3/2}2^{-\ell/2})$, 对几何级数求和得到 $O(N^{3/2})$.

当 $d=3$, 层存储为 $O(N^{4/3}2^{-\ell/3})$, 层工作为 $O(N^2 2^{-\ell})$, 两个级数分别为 $O(N^{4/3})$ 和 $O(N^2)$.

当 $d=1$, 分隔子和前沿大小为常数, 每层成本为 $O(2^\ell)$, 直到 $2^\ell=O(N)$; 总和为 $O(N)$. 最后还需保存和处理叶子区域, 同样为 $O(N)$. 三角求解沿因子条目作乘加, 因而成本为 $O(\operatorname{nnz}(L))$.
\end{proof}
对于二维模型, 最顶层界面虽然只有 $O(\sqrt N)$ 个未知数, 它的稠密分解已经需 $O(N^{3/2})$; 对三维, 顶层 $O(N^{2/3})$ 的界面需 $O(N^2)$. 这解释了为什么三维问题往往在较小的总未知数规模就遭遇直接法内存和时间瓶颈.

这些是模型量级, 不能用来预测每个具体常数. 多重前沿、超节点、消元树上的并行和稠密块 BLAS 可以显著改善实际速度, 却不消除几何界面导致的主要幂次. 非规则、细长、强连接或近乎稠密的图, 需要重新分析分隔子与前沿大小.

\originalsection{数值主元、符号分析与多右端复用}
稀疏直接求解一般分为三阶段. 第一阶段分析图结构, 选择排列, 预测因子结构、消元树和所需工作区; 第二阶段按当前数值执行分解; 第三阶段对一个或多个右端作三角求解. 若数值改变但非零结构不变, 符号阶段通常可以复用. 若矩阵不变、右端改变, 数值因子也可以复用, 这是直接法的重要优势.

SPD Cholesky 不需要为了避免小或负主元而作普通部分选主元, 所以符号结构较容易预先确定. 但浮点下接近奇异、缩放很差的 SPD 输入仍可能暴露数值问题, 不能把结构分析当成准确性证明.

一般非对称稀疏 LU 还要协调两个目标: 结构排列减少填充, 数值主元保证稳定. 可能先采用列排序减少预期填充, 再在分解中进行行主元选择或延迟消元. 数值主元变化会影响实际填充和前沿大小, 所以仅凭符号排序不能完全固定全部成本. 为了更稀疏而接受过小主元, 会增大消元乘子、元素增长和舍入误差; 为了更稳定而改变顺序, 可能增加填充. 二者必须共同判断.

若需要高精度残差或修正解, 可以利用已经得到的因子作迭代改进: 计算真实残差, 用因子近似求解校正方程, 再更新解. 它是否有效仍取决于因子误差、条件数、残差计算精度和舍入范围, 不是任意不稳定因子的自动补救.

\originalsection{直接法和迭代法怎样比较}
一个有用的比较不应只数原矩阵的非零元. 对直接法要估计\emph{填充后的因子}, 对迭代法要估计\emph{达到指定精度的迭代数及预条件成本}. 相同的 $N$ 和 $\operatorname{nnz}(A)$, 可以对应截然不同的计算难度.

以 $d$ 维 Poisson 方程的标准最近邻差分为例, 假设等距网格并采用齐次 Dirichlet 边界. 对每个方向 $m$ 个内部点, 不带 $h^{-2}$ 缩放的离散负 Laplace 特征值为
\[
\lambda_{k_1,\ldots,k_d}
=4\sum_{j=1}^d\sin^2\left(\frac{\pi k_j}{2(m+1)}\right),
\qquad 1\le k_j\le m.
\]
其谱公式可直接推导. 一维取向量 $v_i=\sin(\pi ki/(m+1))$, 边界 $v_0=v_{m+1}=0$. 余弦加法给 $2v_i-v_{i-1}-v_{i+1}=2(1-\cos(\pi k/(m+1)))v_i=4\sin^2(\pi k/(2(m+1)))v_i$. 不同 $k$ 给出一维全部正交特征方向. $d$ 维算子是各坐标的一维算子的 Kronecker 和, 各方向正弦向量的张量积因此是特征向量, 特征值为各一维特征值之和, 得上述公式.

最小值为 $\Theta(m^{-2})$, 最大值为 $\Theta(1)$, 所以 $\kappa_2(A)=\Theta(m^2)=\Theta(N^{2/d})$. 若乘以 $h^{-2}$, 两端同时缩放, 条件数不变. 因而未预条件 CG 的 Chebyshev 保证约需
$O(N^{1/d}\log(1/\tau))$ 步; 每步矩阵向量积为 $O(N)$, 总工作上界约为
\[
O(N^{1+1/d}\log(1/\tau)),\qquad\text{向量存储 }O(N).
\]
二维中这个幂次与嵌套分割分解的 $N^{3/2}$ 相近, 但精度、常数、存储和多右端复用不同. 三维未预条件 CG 的该界为 $N^{4/3}$ 量级, 小于直接分解的 $N^2$ 幂次, 可仍然随网格加密增加迭代数. 若有合适的网格无关多重网格预条件, 迭代数还可能保持有界, 总工作接近线性; 这依赖实际方程和预条件器的有效性, 不能只由稀疏性推出.

选择方法时, 可以按以下问题逐项判断:
\begin{enumerate}
\item 是否只有快速 $Av$, 却没有可访问的矩阵条目? 若是, 普通稀疏直接分解无从实施, Krylov 方法更自然.
\item 是否 SPD、Hermitian 不定或一般非对称? 这决定 CG、MINRES、GMRES 等方法的适用条件.
\item 因子填充能否装入内存? 三维界面、图中高度连接区域和主元变动可能使输入很小而因子很大.
\item 有多少个右端, 矩阵多久改变一次? 反复复用分解可能摊薄直接法的构造成本; 有效预条件器同样可以复用.
\item 所需精度多高, 系统有多病态? 迭代法可能后期停滞, 直接法也可能需要缩放、改进和残差验算.
\item 在并行机器上, 三角求解依赖和全局内积通信各自有多昂贵? 运算量相同不代表实际时间相同.
\end{enumerate}
较小稠密问题优先采用稳定的成熟稠密分解; 因子仍可容纳的稀疏问题常适合稀疏直接法; 最大规模或仅有矩阵向量积的问题常需要预条件迭代. 这些是由结构和资源支持的选择原则, 不应把某个固定维数阈值当作所有硬件和问题通用的定理. 对边界情形, 使用同一真实残差标准做小规模试算最有说服力.

\originalsection{习题与解答}
\begin{exercise}
对长度 $N$ 的路径图, 比较沿路径自然顺序消去, 与先消去一个内部顶点的第一步. 为什么三对角直接法本来就有线性复杂度, 嵌套分割不会改变其主要幂次?
\end{exercise}
\begin{solution}
自然顺序每次消去路径端点, 只剩一个邻居, 没有新边. 先消去内部顶点会连接其两个邻居, 产生一条填充边, 但剩余图仍是路径. 无论适当的递归顺序如何组织, 一维分隔子和前沿都是常数大小, 总成本为 $O(N)$. 排序可能改变常数和并行性, 不能把已是线性的工作降为次线性, 因为仅读写输入输出就需线性成本.
\end{solution}
\begin{exercise}
一个 $N$ 阶 SPD 矩阵有半带宽 $w=N^{1/3}$. 给出带状 Cholesky 的存储和分解上界. 仅凭这个半带宽, 能否断言嵌套分割一定获得二维模型的成本?
\end{exercise}
\begin{solution}
存储为 $O(N^{4/3})$, 分解为 $O(N^{5/3})$. 半带宽只说明当前编号下边的跨度, 不保证存在二维网格式 $O(\sqrt N)$ 分隔子, 也不控制递归区域的前沿. 所以不能直接使用二维嵌套分割公式.
\end{solution}
\begin{exercise}
设两个互不相连的子图通过分隔子 $S$ 连接. 证明先消去两个子图内部时, 一个子图的消元不会在另一个子图的内部顶点之间直接引入由前者产生的新边.
\end{exercise}
\begin{solution}
两个子图的内部之间原来没有边. 前者内部某顶点的剩余邻居只能位于同一内部或 $S$. 把这些邻居连接成团后, 仍没有后者内部顶点参与. 归纳所有内部消元步骤即可. 但 $S$ 上会积累更新, 最终处理 $S$ 时又会产生稠密界面耦合, 这正是分隔子成本必须计入的原因.
\end{solution}
\begin{exercise}
一个稀疏矩阵固定不变, 要求解 $1000$ 个不同右端. 为什么一次直接分解可能优于逐个独立运行 Krylov 方法? 又为什么这个结论不能仅凭右端个数就确定?
\end{exercise}
\begin{solution}
直接法只分解一次, 后续每个右端只做两次三角求解. 独立 Krylov 求解重复生成空间和执行内积, 成本可能很大. 但若因子填充过多装不下, 直接法不可行; 若预条件迭代很快, 或能复用子空间、使用块方法, 迭代总成本也可能更低. 需要比较构造成本、每右端成本、内存与目标精度.
\end{solution}
\begin{exercise}
二维规则网格的每层嵌套分割存储为何都是 $O(N)$, 而每层分解工作却按几何级数递减? 用此解释 $O(N\log N)$ 与 $O(N^{3/2})$ 的不同来源.
\end{exercise}
\begin{solution}
第 $\ell$ 层有 $2^\ell$ 个区域, 分隔子大小约 $(N/2^\ell)^{1/2}$. 平方后乘区域数为 $N$, 每层都贡献同阶存储, 层数 $O(\log N)$. 立方后乘区域数为 $N^{3/2}2^{-\ell/2}$, 越深层越便宜, 所以工作求和由顶层主导, 不再多一个对数因子.
\end{solution}
